\chapter{Food Intolerance and Testing}

\section*{Definition and Overview}
Food intolerance is a reproducible, unpleasant reaction to a food that does not involve the immune system, unlike a food allergy. It is much more common than true food allergy and causes symptoms such as bloating, wind, abdominal pain, diarrhoea, and headache, which usually develop hours after eating the food and are dose-related: small amounts may be tolerated, while larger amounts cause symptoms. The most common intolerances are to lactose, fructose, and other fermentable carbohydrates (FODMAPs), and to food chemicals such as amines, salicylates, and glutamates. Diagnosis is made by careful dietary review and elimination, not by unreliable commercial tests. This chapter explains food intolerance, how it is tested properly, and which tests to avoid.

\section*{Epidemiology}
Food intolerance is common. Lactose intolerance affects around 1 in 20 Australians of European descent but the majority of people of Asian, African, and Indigenous Australian ancestry, because lactase production declines after childhood in most populations. Irritable bowel syndrome, in which food-related symptoms are very common, affects 10--15% of the population, and many of these people find FODMAP foods trigger symptoms. Despite its frequency, food intolerance is often misdiagnosed, and many people spend money on unvalidated "food intolerance tests" such as IgG panels, hair analysis, and kinesiology, which have no scientific basis.

\section*{Aetiology and Pathophysiology}
\begin{itemize}
    \item \textbf{Lactose intolerance:} low levels of the enzyme lactase mean lactose in milk is not digested and is fermented by gut bacteria, causing bloating, wind, and diarrhoea
    \item \textbf{Fructose intolerance:} impaired absorption of fructose from fruit, honey, and some sweeteners causes similar symptoms
    \item \textbf{FODMAPs:} fermentable carbohydrates (fructose, lactose, polyols, fructans, and galactans) are poorly absorbed and fermented, triggering symptoms in sensitive people, particularly those with irritable bowel syndrome
    \item \textbf{Food chemicals:} amines, salicylates, and glutamates, found in many foods including tomatoes, cheese, and additives, cause reactions in sensitive individuals
    \item \textbf{Histamine intolerance:} impaired breakdown of histamine in foods such as aged cheese, wine, and fish causes flushing, headache, and gut symptoms
    \item \textbf{Mechanism:} most intolerance is about how food is digested and absorbed, not immunity
    \item \textbf{Non-IgE mechanisms:} some reactions involve the immune system but not the allergy antibody IgE, such as in coeliac disease, which is an autoimmune reaction to gluten
    \item \textbf{The placebo and nocebo effect:} symptoms can be provoked by expectation, which is why testing must be blinded to be reliable
\end{itemize}

\section*{Clinical Features}
\begin{itemize}
    \item Bloating, wind, and abdominal discomfort, often worse after meals
    \item Diarrhoea, loose stools, or urgent bowel movements
    \item Nausea and, less commonly, vomiting
    \item Headache, fatigue, and a feeling of being unwell
    \item Skin flushing or itching in some chemical and histamine reactions
    \item Symptoms that start hours after eating and are dose-related
    \item The absence of immediate allergic features: no hives, lip swelling, or breathing difficulty
    \item Symptoms that overlap with irritable bowel syndrome, which is the most common context
\end{itemize}

\section*{Differential Diagnosis}
Food-related symptoms must be distinguished from other conditions.
\begin{itemize}
    \item Food allergy, which is immune-mediated and can be immediate and severe
    \item Coeliac disease, which requires testing with a gluten-containing diet
    \item Irritable bowel syndrome, which is the most common underlying condition
    \item Inflammatory bowel disease, which needs assessment for blood, weight loss, and inflammation
    \item Gallbladder and pancreatic conditions causing fat intolerance
    \item Infection and other causes of diarrhoea
    \item The danger of assuming intolerance and missing these other diagnoses
\end{itemize}

\section*{When to Seek Medical Care}
Anyone with persistent food-related symptoms should see a doctor for assessment, because symptoms of intolerance overlap with serious conditions. Medical advice is essential before starting restrictive elimination diets, which can cause nutritional deficiency.

\textbf{Call 000 (or local emergency services) for:}
\begin{itemize}
    \item Difficulty breathing, throat swelling, or collapse after eating, which indicates anaphylaxis, not intolerance
    \item Severe abdominal pain, especially with vomiting
    \item Blood in the stool or black stools
    \item Signs of severe dehydration: no urine, dizziness, or confusion
\end{itemize}

\section*{Diagnosis and Investigations}
\begin{itemize}
    \item \textbf{Clinical history:} the food, the amount, the timing, and the symptoms are the most important diagnostic information
    \item \textbf{Food and symptom diary:} recording intake and symptoms over a few weeks identifies patterns
    \item \textbf{Elimination and challenge:} removing the suspected food for 2--4 weeks, then reintroducing it to confirm the reaction
    \item \textbf{Lactose hydrogen breath test:} measures hydrogen in the breath after a lactose drink; the gold standard for lactose intolerance
    \item \textbf{Other breath tests:} for fructose and small intestinal bacterial overgrowth
    \item \textbf{FODMAP elimination diet:} a structured, dietitian-guided program for irritable bowel syndrome
    \item \textbf{Coeliac testing:} blood tests for coeliac disease must be done before any gluten elimination
    \item \textbf{Tests to avoid:} IgG food panels, hair analysis, kinesiology, and other unvalidated tests, which are not supported by evidence and can mislead
\end{itemize}

\section*{Management}
\begin{itemize}
    \item \textbf{Lactose intolerance:} limit or avoid milk, or use lactose-free dairy and lactase supplements
    \item \textbf{FODMAP intolerance:} a dietitian-guided low-FODMAP diet in three phases: elimination, structured reintroduction to identify triggers, and a personalised maintenance diet
    \item \textbf{Chemical intolerance:} an elimination diet under supervision to identify and then minimise trigger foods, ensuring nutrition is maintained
    \item \textbf{Small amounts:} many people tolerate small servings of trigger foods even when larger servings cause symptoms
    \item \textbf{Treat coexisting conditions:} irritable bowel syndrome responds to dietary, lifestyle, and sometimes medical treatment
    \item \textbf{Nutritional safety:} dietary elimination must be supervised to prevent deficiency, especially in children and pregnancy
    \item \textbf{Supplements:} lactase enzymes for lactose, and alpha-galactosidase for some legumes, help some people
\end{itemize}

\section*{Complications}
\begin{itemize}
    \item Unnecessary, over-restrictive diets causing weight loss and nutrient deficiency
    \item Missed diagnosis of coeliac disease or inflammatory bowel disease
    \item Wasted money on unvalidated testing
    \item Reduced quality of life and social restriction from fear of foods
    \item Anxiety about eating and food avoidance in children when intolerance is over-diagnosed
\end{itemize}

\section*{Prognosis and Outcomes}
Most food intolerances are well managed once the trigger is identified, and symptoms settle with dietary adjustment. Lactose intolerance is lifelong but easily handled with lactose-free options, and most people with FODMAP sensitivity find their personal triggers through structured reintroduction and return to a broad diet. The key to a good outcome is proper diagnosis: a validated approach identifies true triggers, while commercial testing leads to confusion and unnecessary restriction. With dietitian guidance, people with food intolerance eat well and live normally.

\section*{Prevention and Self-Care}
\begin{itemize}
    \item See a doctor first to exclude conditions such as coeliac disease and inflammatory bowel disease
    \item Keep a food and symptom diary before eliminating anything
    \item Trial elimination of one food at a time for 2--4 weeks, then reintroduce
    \item Use breath testing for lactose and fructose where available
    \item Work with a dietitian for FODMAP and other elimination diets
    \item Avoid unvalidated commercial intolerance tests
    \item Do not put a child on a restrictive diet without medical advice
    \item Reintroduce foods gradually to maintain the widest possible diet
\end{itemize}

\section*{Sources and Further Reading}
The following reputable sources provide further information for healthcare students and patients seeking reliable, current advice about food intolerance and testing.
\begin{itemize}
    \item Healthdirect Australia. \textit{Food intolerance.} https://www.healthdirect.gov.au/food-intolerance
    \item Dietitians Australia. \textit{Food intolerance and FODMAPs.} https://dietitiansaustralia.org.au/
    \item The Gut Foundation. \textit{Food intolerance.} https://www.gutfoundation.org.au/
    \item Monash University. \textit{Low FODMAP diet.} https://www.monashfodmap.com/
    \item NHS. \textit{Food intolerance.} https://www.nhs.uk/conditions/food-intolerance/
    \item World Health Organization. \textit{Food safety.} https://www.who.int/
\end{itemize}
